\section{Getting Rich, Making Something, and Becoming Someone}

\subsection{What the Business Produces in the Founder}

There is 1 last distinction to make. People sometimes imagine a career as a sequence of external achievements: a higher salary, a more impressive title, a larger company, an investment round, a valuation, a famous office, or a campus. Such markers can provide useful information. They are also dangerously easy to optimize independently of the underlying goal.

A founder who wants to become wealthy may be better served by a boring, profitable company than a fashionable startup. A founder who wants to solve an important scientific problem may need to spend years building credibility rather than revenue. A founder who wants to create a great school may need to refuse opportunities that would produce short-term enrollment at the cost of learning quality. A founder who wants autonomy may need to avoid funding structures that require a particular kind of scale.

To understand a business, ask what it is producing. To understand the founder, ask what the business is producing in the founder.

Does it develop better judgment? Does it create the freedom the founder hoped for, or more dependence? Does it make them more intellectually honest, or more attached to defending an old hypothesis? Does it develop the capacity to work with extraordinary people? Does it make promises they are proud to keep? Does it gradually create assets and relationships that make the next generation of work possible?

\textbf{A business that earns money while slowly destroying the founder's curiosity, relationships, or integrity may have an impressive financial statement and a poor total return.} Conversely, a company that fails commercially can still create valuable knowledge and relationships, although those benefits should not be invoked as a retrospective justification for taking reckless risks. Not every failure is a disguised success.

The purpose of these questions is not to moralize about wanting money. It is to recognize that money is usually sought because of what it makes possible. If a founder forgets that, wealth can become another treadmill rather than an escape from 1.
